% ============================================================================
\section{Hasil dan pembahasan}
\label{sec:hasil}
% ============================================================================

\isian{Bagian ini diisi setelah E0--E4 dijalankan sesuai pra-registrasi. Untuk setiap RQ dilaporkan metrik utama dengan CI 95\% dan putusan hipotesisnya mengikuti Tabel~\ref{tab:lacak}, hasil per jenis perubahan dan per lapisan, perbandingan dengan pembanding dan ablasi, serta analisis kesalahan LLM menurut taksonomi yang ditetapkan sebelum pelabelan.}

\subsection{Ancaman terhadap validitas}
Acuan dan \en{engine} disusun tim yang sama, sehingga kesalahan yang sama dapat muncul di keduanya. Hal ini dimitigasi dengan acuan pihak ketiga, subset \en{holdout} contoh resmi, dan pengadjudikasi kedua. Representasi dirancang setelah perubahan 2023--2026 diketahui, sehingga klaim adaptivitas juga diuji pada perubahan tak terlihat dan peristiwa prospektif. Kecocokan dengan acuan belum tentu berarti kesesuaian hukum, karena penelitian ini tidak melibatkan pakar hukum pajak. Hasilnya terbatas pada satu yurisdiksi, satu jenis pajak, dan satu rangkaian perubahan dengan jumlah peristiwa yang kecil. Regulasi 2023--2026 juga mungkin ada di data latih LLM, dan hal ini diperiksa dengan dokumen kontrafaktual dan perbandingan hasil sebelum dan sesudah \en{training cutoff} model.
